(1) 令 \(x(t)=\Re[Xe^{-i\omega t}]\)，代入有 \((\omega_0^2-\omega^2-2i\gamma\omega)X=F_0/m\)，所以
\[
|X|=\frac{F_0/m}{\sqrt{(\omega_0^2-\omega^2)^2+4\gamma^2\omega^2}}.
\]

(2) 要最大化位移幅度，即最小化分母平方 \(D=(\omega_0^2-\omega^2)^2+4\gamma^2\omega^2\)。求导 \(D'=4\omega(\omega^2-\omega_0^2+2\gamma^2)\)，因此在 \(\gamma<\omega_0/\sqrt2\) 下正峰位于
\[
\omega_{\rm amp}=\sqrt{\omega_0^2-2\gamma^2}<\omega_0.
\]
阻尼越大，位移峰向低频移动。

(3) 稳态能量由外力输入并全部由阻尼耗散，平均功率
\[
\overline P=2m\gamma\overline{\dot x^2}
=m\gamma\omega^2|X|^2
=\frac{\gamma F_0^2\omega^2/m}{D}.
\]
对 \(\omega^2/D\) 求导，极大值在 \(\omega=\omega_0\)。所以功率共振峰位于固有无阻尼频率，而位移共振峰低于它。

(4) 当 \(\gamma\to0^+\)，两共振频率均趋 \(\omega_0\)，峰值位移 \(|X|_{\rm peak}\propto1/\gamma\) 发散，平均峰值功率 \(\overline P_{\max}=F_0^2/(4m\gamma)\) 也发散。严格 \(\gamma=0\) 时共振响应非稳态周期函数，不能把公式的无限值解释为有限时间内的实际振幅。
